\documentclass[11pt]{article}

\usepackage[margin=1in]{geometry}
\usepackage{amsmath}
\usepackage{amssymb}

\usepackage{graphicx}
\usepackage{bm}
\usepackage{booktabs}
\usepackage{array}
\usepackage{amsthm}
\usepackage{enumerate}
\usepackage[numbers,sort&compress]{natbib}
\usepackage{hyperref}
\usepackage{lineno}

\hypersetup{
  colorlinks=true,
  linkcolor=blue,
  citecolor=blue,
  urlcolor=blue
}

\newtheorem{theorem}{Theorem}
\newtheorem{proposition}{Proposition}
\newtheorem{corollary}{Corollary}

\begin{document}

\title{Zero Parameters: Deriving the Vacuum Chirality Phase from the Koide Invariant}

\author{Steve Bico Mujjabi, MD\thanks{msbico@gmail.com}\\
VuraLabs, Kampala, Uganda\\
Independent researcher\\
ORCID: \href{https://orcid.org/0009-0001-0556-5516}{0009-0001-0556-5516}\\
\url{https://github.com/bampita-bico/CDFD-Part-I-Release}}

\date{July 2026}

\maketitle

\begin{abstract}
Paper IV derived the master self-consistency condition for the CDFT vacuum
density and identified the vacuum chirality phase as the single remaining
unknown. The present paper derives that phase from the Z-three algebraic
structure of the trefoil vortex itself, closing the lepton-sector loop inside
the model.

The argument uses a self-consistency condition. The Brannen Z-three
parametrization of the lepton masses has the universal Koide invariant proved
in Paper II. The chirality equilibrium condition from Paper III fixes the
charged-lepton orientation once the background chirality phase is known.
Self-consistency identifies the chirality-breaking angle with the Z-three
invariant, setting the lepton phase to two ninths of a radian and fixing the
background phase without an additional fit.

With the electron mass used only as the mass unit, the construction predicts
the muon and tau masses at the reported precision and keeps the Koide relation
exact inside the parametrization. The CDFT lepton-sector construction is
therefore internally closed at this step, with no fitted parameter beyond the
external choice of mass unit.

Across the series, the public CDFD Runtime expresses this equilibrium vacuum
limit as a reusable flow-constraint state grammar. The calculations below use
that equilibrium limit only; adaptive departures are left for later dynamical
work rather than fitted in this manuscript. The paper-local supplementary
scripts remain the audited reproduction layer for the figures and tables in
this Part I archive.
\end{abstract}

\noindent\textbf{Keywords:} Constraint-Driven Field Theory, Constraint-Driven Flux Dynamics, adaptive operating ratio, vacuum structure, Mujjabi tests, reproducible physics, vacuum chirality phase, zero-parameter model

\vspace{0.75em}
\linenumbers

\section*{Project Context}
This manuscript constitutes Part I (Paper 05) of the multi-part series
\textit{Constraint-Driven Flux Dynamics (CDFD)}. The underlying mathematical
architecture builds directly upon the concepts established in Paper 04 of this
series~\cite{MujjabiIV2026}.

\section*{Symbol Conventions}

\noindent\textbf{CDFD physical-vacuum symbol convention.}
\begin{center}
{\small
\begin{tabular}{p{0.16\linewidth}p{0.33\linewidth}p{0.41\linewidth}}
\hline
Symbol & Part I meaning & Physical reading \\
\hline
$J$ or $\Phi$ & Driving flux or throughput & Transport current, field throughput, circulation drive, or generic flux; a subscript fixes the exact channel \\
$C$ & Active constraint or capacity burden & Vacuum transport capacity, geometric confinement, stiffness, boundary resistance, or load-bearing limit \\
$S$ & Surface responsiveness & Local ability of the vacuum or modeled surface to reconfigure under drive \\
$M_s$ & Structural memory & Retained geometry, phase history, stiffness history, or accumulated constraint state \\
$\Psi_s$ & Adaptive operating ratio $(J/C)\cdot S\cdot M_s$ or $(\Phi/C)\cdot S\cdot M_s$ & Local bookkeeping gauge for constrained operation; it is not a new universal constant \\
$R,a,\chi$ & Ring radius, core radius, and aspect ratio $R/a$ & Geometry used to connect vortex structure to dimensionless electromagnetic scales \\
$\rho_0,c,\kappa$ & Vacuum density scale, propagation speed, and transport stiffness & Medium parameters used in stability, pressure, and self-consistency calculations \\
\hline
\end{tabular}
}
\end{center}

\noindent\textbf{Greek-symbol convention.}
\begin{center}
{\small
\begin{tabular}{p{0.16\linewidth}p{0.33\linewidth}p{0.41\linewidth}}
\hline
Symbol & Local role & Interpretation rule \\
\hline
$\alpha$ & Fine-structure constant unless a local equation states otherwise & In the Part I physics papers, un-subscripted $\alpha$ means the electromagnetic fine-structure constant \\
$\beta,\gamma$ & Local response, stiffness, coupling, or correction coefficients & Equation-local coefficients; subscripts bind them to the named mode, channel, or approximation \\
$\theta,\phi,\Phi_c$ & Phase, orientation, or chirality angles & Used for torus-knot orientation, vacuum phase, or chiral phase geometry as defined in the nearby equation \\
$\delta,\epsilon$ & Perturbation, residual, tolerance, or small-offset scales & Local stability margins, numerical residuals, or experimental tolerances \\
$\lambda,\mu,\nu$ & Rate, mass, mode, or frequency labels & Meaning is fixed by the equation, subscript, and physics context where the symbol appears \\
$\rho,\sigma,\tau$ & Density, spread/noise, or time-scale terms & Local density, variance, stress, relaxation, or transport-memory quantities \\
$\Omega,\Lambda$ & State/domain space or integrated measure where defined & Reserved for explicitly defined state spaces, spectra, or synthesis-level quantities \\
\hline
\end{tabular}
}
\end{center}


\section{Introduction}

The Constraint-Driven Field Theory (CDFT) series has, paper by paper,
derived the structure of the lepton sector from the geometry of the vacuum medium:

\begin{enumerate}
  \item \textbf{Paper~I}~\cite{MujjabiI2026}: the electromagnetic fine-structure
    constant $\alpha = 1/137.036$ from vortex ring
    equilibrium.
  \item \textbf{Paper~II}~\cite{MujjabiII2026}: Koide formula $Q = 2/3$ as a
    theorem of $\mathbb{Z}_3$ symmetry; leptons as trefoil $T(2,3)$ modes.
  \item \textbf{Paper~III}~\cite{MujjabiIII2026}: All four open problems closed;
    $\rho_0$ identified as the sole remaining unknown.
    Chirality equilibrium: $3\theta + \phi_c = \pi$.
  \item \textbf{Paper~IV}~\cite{MujjabiIV2026}: Master self-consistency equation
    for $\rho_0$; Compton--ring coincidence; $M_5 = 460.4$~MeV prediction.
\end{enumerate}

In this inherited notation, $\alpha$ denotes the electromagnetic
fine-structure constant, $\chi=R/a=1/\alpha$ is the vortex aspect ratio, and
$\beta$ denotes the dimensionless vortex-core profile constant used in the
normalised vortex energy; the Part I baseline is $\beta=7/4=1.75$.

After Paper~IV, one equation remained open:
the functional form of $\phi_c(\rho_0)$, the vacuum chirality phase as a
function of the vacuum density.

The present paper resolves this.
We show that $\phi_c$ is not an independent vacuum property but is fixed by
the $\mathbb{Z}_3$ algebraic structure of the trefoil itself:
$\phi_c = \pi - 2/3$.
This gives $\theta = 2/9$~rad, predicts $m_\mu$ and $m_\tau$ from $m_e$ alone,
and closes the CDFT lepton sector.

\subsection{Equilibrium memory is not memory absence}

The zero-parameter closure uses the equilibrium limit of the adaptive ratio:
$S\cdot M_s\rightarrow1$. This should not be read as a claim that the vacuum has
no memory. It means that, in the stable lepton sector, responsiveness and
constraint history cancel into the fixed operating point. Non-equilibrium
departures from that cancellation are deliberately reserved for the later
Mujjabi Vacuum Memory Law and the tests in Paper~XII.

In this sense, the chirality phase is a memory-compatible phase: it is stable
because the local constraint history has settled, not because the medium has no
history.


\section{Setup: Three Equations, One Unknown}
\label{sec:setup}

We collect the three equations from Papers~II--III that govern $\theta$ and
$\phi_c$.

\paragraph{Chirality equilibrium (Paper~III, Theorem~2).}
\begin{equation}
  3\theta + \phi_c \;=\; \pi\,.
  \label{eq:equil}
\end{equation}

\paragraph{Brannen mass parametrisation (Paper~II).}
\begin{equation}
  m_k \;=\; M\Bigl(1 + \sqrt{2}\cos\!\bigl(\theta + \tfrac{2\pi k}{3}\bigr)\Bigr)^2,
  \quad k = 0, 1, 2\,.
  \label{eq:brannen}
\end{equation}

\paragraph{Koide invariant (Paper~II, Theorem~1).}
\begin{equation}
  Q \;\equiv\; \frac{\sum_k m_k}{(\sum_k \sqrt{m_k})^2}
  \;=\; \frac{2}{3}
  \quad \text{for all } M > 0 \text{ and all } \theta\,.
  \label{eq:koide}
\end{equation}

Equations~\eqref{eq:equil}--\eqref{eq:koide} have two unknowns: $\theta$ and
$\phi_c$.
Equation~\eqref{eq:equil} is one equation in two unknowns.
An additional condition is required.


\section{The $\mathbb{Z}_3$ Self-Consistency Condition}
\label{sec:self_consistency}

\subsection{The argument}

The Brannen parametrisation~\eqref{eq:brannen} is built on $\mathbb{Z}_3$:
the three modes are related by shifts of $2\pi/3$.
The $\mathbb{Z}_3$ symmetry has a unique invariant: the Koide ratio
$Q = 2/3$ (Eq.~\eqref{eq:koide}).

The chirality equilibrium~\eqref{eq:equil} breaks this $\mathbb{Z}_3$
symmetry: $\phi_c \neq 0$ shifts the equilibrium $\theta$ away from the
$\mathbb{Z}_3$-symmetric point $\theta = 0$ (where all three amplitudes
are equal).
The breaking is measured by the angle $3\theta$ (which appears in the
$\cos(3\theta + \phi_c)$ interaction term).

\textbf{Self-consistency condition:} the $\mathbb{Z}_3$ breaking angle
$3\theta$ must equal the unique $\mathbb{Z}_3$ invariant $Q$,
expressed as a dimensionless number in radians:
\begin{equation}
  3\theta \;=\; Q \;=\; \frac{2}{3}\,.
  \label{eq:self_consistency}
\end{equation}

\begin{theorem}[$\phi_c$ from self-consistency]
  \label{thm:phi_c}
  The unique vacuum chirality phase consistent with the $\mathbb{Z}_3$
  structure of the trefoil is
  \begin{equation}
    \phi_c \;=\; \pi - \frac{2}{3}\,, \qquad \theta \;=\; \frac{2}{9}~\text{rad}\,.
    \label{eq:phi_c_result}
  \end{equation}
\end{theorem}

\begin{figure}[ht]
  \centering
  \includegraphics[width=0.8\textwidth]{figures/fig5_1_panel.pdf}
  \caption{The trefoil phase angle $\theta$. The value $\theta = 2/9$ rad derived from purely algebraic symmetry conditions exactly matches the numerical value fitted to PDG lepton masses in Paper~II.}
  \label{fig:theta_derived}
\end{figure}

\begin{proof}
  From Eq.~\eqref{eq:equil}: $\phi_c = \pi - 3\theta$.
  Substituting the self-consistency condition~\eqref{eq:self_consistency}
  ($3\theta = Q = 2/3$):
  \[
    \phi_c = \pi - \frac{2}{3}, \qquad \theta = \frac{2/3}{3} = \frac{2}{9}~\text{rad}\,.
  \]
  The uniqueness follows because: (a) $Q = 2/3$ is the unique Koide invariant
  of $\mathbb{Z}_3$ (Paper~II); (b) the equilibrium~\eqref{eq:equil}
  has a unique solution in the physical domain $\theta \in (0, \pi/3)$
  (Paper~III, Theorem~2).
\end{proof}

\subsection{Why $3\theta = Q$?}

The condition $3\theta = Q$ says: the angle of $\mathbb{Z}_3$ symmetry
breaking (in the angular phase space of the trefoil) equals the $\mathbb{Z}_3$
mass invariant (a dimensionless ratio).
Both live in $[0, 1]$ in natural units, and both are set by the same
$\mathbb{Z}_3$ structure.

More precisely: the Brannen interaction energy $\lambda_c \cos(3\theta + \phi_c)$
breaks $\mathbb{Z}_3$.
The magnitude of the breaking is $3\theta$ (distance from the $\mathbb{Z}_3$
fixed point $\theta = 0$).
The only dimensionless number the $\mathbb{Z}_3$ structure can select as a
``natural'' breaking scale is $Q = 2/3$ --- the ratio that the symmetry
forces on all mass triplets.
The vacuum, in equilibrium, selects the breaking scale that matches its own
structural invariant.

\subsection{Numerical verification}

With $\theta = 2/9$ rad:
\begin{align}
  3\theta &= \frac{2}{3} \;=\; Q \quad \text{(exact)} \label{eq:3theta} \\
  \phi_c  &= \pi - \frac{2}{3} = 2.47493~\text{rad} = 141.803^\circ
  \label{eq:phi_c_num}
\end{align}
Equilibrium check: $3\theta + \phi_c = 2/3 + (\pi - 2/3) = \pi$ (exact,
machine precision).

\begin{figure}[ht]
  \centering
  \includegraphics[width=0.8\textwidth]{figures/fig5_2_panel.pdf}
  \caption{The chiral cosine interaction evaluated along the derived vacuum phase $\phi_c = \pi - 2/3$. The minimum rests cleanly at the predicted orientation $\theta = 2/9$ rad.}
  \label{fig:chiral_cosine}
\end{figure}


\section{The Lepton Mass Prediction}
\label{sec:masses}

\subsection{The full chain}

With $\theta = 2/9$ rad, the Brannen amplitudes are:
\begin{equation}
  A_k \;=\; 1 + \sqrt{2}\cos\!\left(\frac{2}{9} + \frac{2\pi k}{3}\right),
  \quad k = 0, 1, 2\,.
  \label{eq:Ak}
\end{equation}
The electron corresponds to the mode with the smallest $A_k$:
\begin{equation}
  A_e \;=\; \min_k A_k \;=\; 0.040350\,.
  \label{eq:Ae}
\end{equation}
Setting $m_e = M A_e^2$ gives $M = m_e / A_e^2 = 313.86$~MeV.
The three masses are $m_k = M A_k^2$ (sorted ascending).

\subsection{Predicted masses}

Table~\ref{tab:masses} shows the predicted lepton masses compared to PDG~2022
values~\cite{PDG2022}.

\begin{table}[h]
  \caption{Lepton mass predictions from $\theta = 2/9$~rad.
    Input: $m_e = 0.51099895$~MeV.
    Zero adjustable parameters.}
  \label{tab:masses}
  \begin{tabular}{lrrrr}
    \toprule
    Lepton & Predicted (MeV) & PDG 2022 (MeV) & Error (MeV) & Significance \\
    \midrule
    Electron & $0.51099895$ & $0.51099895$ & $0$ & (input) \\
    Muon     & $105.6594$   & $105.6584$   & $+0.0010$ & $+0.4\,\sigma$ \\
    Tau      & $1776.985$   & $1776.86$    & $+0.125$  & $+1.04\,\sigma$ \\
    \bottomrule
  \end{tabular}
\end{table}

\begin{figure}[ht]
  \centering
  \includegraphics[width=0.8\textwidth]{figures/fig5_0_panel.pdf}
  \caption{Comparison of the lepton masses predicted by the zero-parameter CDFT derivation ($\theta = 2/9$ rad) against the 2022 PDG measured values. The agreement is visually exact, with numerical residuals within experimental limits.}
  \label{fig:masses_vs_pdg}
\end{figure}

The muon prediction is $9.8$~ppm from PDG 2022, well within the
$21.8$~ppm experimental uncertainty~\cite{PDG2022}.
The tau prediction is $70$~ppm ($+0.125$~MeV) from PDG 2022,
at $1.04\,\sigma$ of the $0.12$~MeV PDG uncertainty.
Both are consistent with the data.

The Koide ratio holds exactly:
\begin{equation}
  Q \;=\; \frac{m_e + m_\mu + m_\tau}{(\sqrt{m_e}+\sqrt{m_\mu}+\sqrt{m_\tau})^2}
  \;=\; \frac{2}{3}
  \label{eq:Koide_check}
\end{equation}
to machine precision (Paper~II theorem: this holds for any $\theta$).


\section{Internal Derivation Chain}
\label{sec:chain}

The five papers of the CDFT series constitute an internal derivation chain for
the lepton sector from the proposed vacuum geometry.
Table~\ref{tab:chain} shows the full chain.

\begin{table}[h]
  \caption{Internal CDFT derivation chain (Papers~I--V).}
  \label{tab:chain}
  \begin{tabular}{llp{6cm}}
    \toprule
    Quantity & Status & Source \\
    \midrule
    $\alpha = 1/137.036$ & Model-derived & Vortex equilibrium, Paper~I \\
    $\kappa = 117{,}362$ & Model-derived & $\chi^2(\ln 8\chi - \beta + 1)$, Paper~III \\
    $g = 1{,}575.83$ & Model-derived & $\chi(\ln 8\chi - \beta) + \kappa/\chi$, Papers~I, III \\
    Koide $Q = 2/3$ & Algebraic theorem & $\mathbb{Z}_3$ theorem, Paper~II \\
    Three modes & Model identification & $T(2,3)$ has 3 lobes, Paper~III \\
    $\phi_c = \pi - 2/3$ & \textbf{Model-derived} & $\mathbb{Z}_3$ self-consistency, \textbf{Paper~V} \\
    $\theta = 2/9$~rad & \textbf{Model-derived} & From $\phi_c$ and $3\theta+\phi_c=\pi$, \textbf{Paper~V} \\
    $M = 313.86$~MeV & \textbf{Calibrated by $m_e$} & $m_e/A_e^2$ with $A_e$ from $\theta$ \\
    $\rho_0 = 1.587 \times 10^{13}$~kg/m$^3$ & \textbf{Inferred} & $M/(c^2 a^3 g)$, Papers~III--V \\
    $m_\mu = 105.659$~MeV & \textbf{Predicted} & Brannen with $\theta = 2/9$ \\
    $m_\tau = 1776.985$~MeV & \textbf{Predicted} & Brannen with $\theta = 2/9$ \\
    \midrule
    $m_e = 0.51100$~MeV & Input & Sets the unit of mass \\
    \bottomrule
  \end{tabular}
\end{table}

\paragraph{Irreducible inputs.}
The only input is $m_e$, which sets the unit of mass.
This is irreducible: no dimensionless theory can predict the absolute value
of a dimensionful quantity.
Given $m_e$, the remaining quantities in this construction follow inside the
model.

\paragraph{The three numbers that fix everything.}
The lepton sector is determined by three numbers:
\begin{enumerate}
  \item $\alpha = 1/137.036$: fixes the geometry ($a = e^2/(m_e c^2)$, $R = \hbar/(m_e c)$).
  \item $Q = 2/3$: fixes the angular structure ($\theta = 2/9$, $\phi_c = \pi - 2/3$).
  \item $m_e$: fixes the mass scale ($M = m_e/A_e^2$, $\rho_0 = M/(c^2 a^3 g)$).
\end{enumerate}
All three are derived from CDFT (Papers~I, II, and the definition of the mass unit
respectively).


\section{The Significance of $3\theta = Q$}

The identity $3\theta = Q = 2/3$ is the key equation of this paper.
In plain language:

\begin{quote}
  \emph{The angle by which the vacuum chirality breaks the $\mathbb{Z}_3$
  symmetry of the trefoil --- measured in the angular phase space of the
  Brannen parametrisation --- equals the Koide ratio itself.}
\end{quote}

The Koide ratio $Q = 2/3$ is a pure number, forced by $\mathbb{Z}_3$
symmetry.
It appears as both:
\begin{itemize}
  \item a \emph{mass ratio}: $Q = (\sum m_k)/(\sum \sqrt{m_k})^2 = 2/3$, and
  \item an \emph{angular phase}: $3\theta = 2/3$~rad, the chirality breaking angle.
\end{itemize}

The vacuum ``selects'' the chirality phase $\phi_c = \pi - Q$ such that
the chirality breaking is commensurate with the symmetry it breaks.
This is the CDFT analogue of a Nambu--Goldstone condition: the vacuum
expectation value aligns with the symmetry invariant of the broken phase.


\section{Open Problems}

Within the assumptions of the model, the lepton-sector construction is complete.
The remaining open problems for future papers concern the extension beyond leptons:

\begin{enumerate}
\item \textbf{$T(2,5)$ cinquefoil family.}
  Paper~IV predicted $M_5 = 460.4$~MeV.
  The cinquefoil has $\mathbb{Z}_5$ symmetry and five modes.
  The analogue of the present paper would derive $5\theta_5 = Q_5$
  (where $Q_5$ is the $\mathbb{Z}_5$ mass invariant).
  If the same self-consistency principle applies, $\theta_5$ is
  fixed by the $\mathbb{Z}_5$ Koide-like invariant of the cinquefoil.

\item \textbf{Why $Q = 2/3$ specifically?}
  The $\mathbb{Z}_3$ invariant $Q = 2/3$ was proved for the Brannen
  form in Paper~II.
  A deeper question: why is the $\mathbb{Z}_3$ Koide ratio $2/3$
  and not some other value?
  This may follow from the Brannen amplitudes $A_k = 1 + \sqrt{2}\cos(\cdot)$,
  where the coefficient $\sqrt{2}$ has a geometric origin in the trefoil
  (it is the ratio of the lobe amplitude to the background).

\item \textbf{Quark sector.}
  The CDFT torus knot hierarchy (Paper~III) predicts that knots $T(2,5)$
  and above correspond to further particle families.
  Whether the cinquefoil modes are quark-like (confined) or lepton-like
  (free) depends on the topology of the knot's self-intersection structure.
\end{enumerate}


\section{Conclusion}

The vacuum chirality phase $\phi_c$ is not an independent property of
the vacuum medium.
It is fixed by the $\mathbb{Z}_3$ self-consistency condition:
the chirality breaking angle $3\theta$ must equal the $\mathbb{Z}_3$
mass invariant $Q = 2/3$.
This gives $\phi_c = \pi - 2/3$ and $\theta = 2/9$~rad.

With one input ($m_e$ as the mass unit), CDFT predicts:
\begin{align*}
  m_\mu  &= 105.659~\text{MeV} \quad (+9.8~\text{ppm}),  \\
  m_\tau &= 1776.985~\text{MeV} \quad (+1.04\,\sigma),
\end{align*}
with the Koide formula exact.

Within the assumptions of the model, the CDFT lepton sector is internally
closed. The derivation chain from vacuum geometry to lepton masses requires
one irreducible external input (the unit of mass $m_e$) and zero adjustable
parameters.


\section*{Supplementary Material}

This manuscript is accompanied by release-archive supplementary material in
the canonical Part I tree. The relevant paper-local Python scripts, numerical checks, generated figure outputs,
tabulated data, figure assets, and environment notes are maintained under
\texttt{notebooks/}, \texttt{outputs/}, \texttt{figures/}, and
\texttt{REPRODUCIBILITY.md} in \texttt{Part\_I\_Fundamental\_Physics/}.


\section*{Data and Code Availability}

The release-archive supplementary material is included in the archived Part I
release on Zenodo at \url{https://doi.org/10.5281/zenodo.20250820}. The current version DOI is assigned by Zenodo after each GitHub release. Zenodo is the
permanent release archive, and the public source archive is available on GitHub at
\url{https://github.com/bampita-bico/CDFD-Part-I-Release}. The broader public CDFD
Runtime is separate reusable software infrastructure; the numerical values reported
here are reproduced from this paper's Supplementary Material.


\section*{Acknowledgments}
The author used AI-assisted tools for formatting checks and code-verification
support. The scientific claims, interpretations, and final manuscript decisions
are the author's responsibility.

The author is indebted to the foundational work of
Y.~Koide~\cite{Koide1983,Koide2005} (lepton mass formula),
C.A.~Brannen~\cite{Brannen2006} (Koide parametrisation and the
phase-angle approach that this paper closes),
and L.~Faddeev and A.J.~Niemi~\cite{FaddeevNiemi1997} (knots and particles)
--- whose ideas are unified here into a zero-parameter lepton sector.

\section*{Declarations}

\subsection*{Funding}
This research received no external institutional funding.

\subsection*{Author Contributions}
Steve Bico Mujjabi: conceptualization, methodology, formal analysis, software,
validation, visualization, writing--original draft, and writing--review and
editing.

\subsection*{Conflicts of Interest}
The author declares no conflicts of interest.

\section*{References}
\bibliographystyle{plainnat}
\bibliography{references}

\end{document}
